Large language model agents now mediate billions of tool calls per day
across customer support, financial services, software engineering, and
clinical workflows. Every one of those calls is the output of a
stochastic system whose input surface is, by construction, the open set
of all natural language. Four years of work on prompt-injection defence
has converged on a stable observation: text-side defences against
text-side attacks improve attack-success rates from catastrophic to
merely poor. The strongest published results in 2024-2026 ---
Spotlighting {[}@hines2024defending{]}, StruQ {[}@chen2024struq{]},
constitutional classifiers {[}@anthropic2024constitutional{]}, and the
suite of commercial prompt-shield products --- reduce the AgentDojo
{[}@debenedetti2024agentdojo{]} and InjecAgent {[}@zhan2024injecagent{]}
attack success rate (ASR) from a no-defence baseline of roughly half to
between 14\% and 31\%. Half of one in three attempts still succeeds.

The diagnosis we offer is that the defences are operating in the wrong
shape of problem. Verifying properties of unbounded natural language is
open; the agent's downstream effect on the world, however, is mediated
almost entirely through a much narrower interface. Every modern agent
framework --- OpenAI function-calling, Anthropic tool use, the Model
Context Protocol, AutoGen, and LangChain among them --- requires tools
to declare structured input grammars, typically JSON Schema. Over those
grammars the set of well-formed strings is bounded, the propositions one
wants to enforce are decidable, and the enforcement point can be moved
out of the model entirely.

This paper presents Verified Capability Discipline (VCD), a composition
of three primitives that together provide \textbf{portable,
cryptographically-verifiable authorisation} for LLM agent tool calls. We
use SMT to prove that a tool's security policy is satisfied over every
string its JSON Schema permits, or to extract a concrete counterexample
call. We bind the proven policy into an Ed25519-signed capability token
whose attenuation primitives mechanically forbid permission broadening.
We place a gate on the only path from agent intent to tool execution and
require every call to carry a token whose constraints are satisfied by
the actual call arguments. The artefact a deployment produces --- the
\emph{capability receipt}, a signed record citing the issuer, the schema
hash, the policy proof hash, the Lean theorem identifier, the
attenuation chain, and a hash of the actual call arguments --- is
portable across organisations, clouds, and trust boundaries, and is the
kind of evidence that regulators, internal auditors, and downstream tool
implementations can independently verify without sharing a secret with
the issuer.

We frame the contribution as authorisation rather than defence
deliberately. The prompt-injection literature has spent five years
asking whether the \emph{model} can be made robust to adversarial input;
in parallel, model-side advances such as instruction/data embedding
separation (ASIDE {[}@zverev2025aside{]}) and preference-aligned
training (SecAlign {[}@chen2025secalign{]}) are pushing baseline ASRs
downward at the model layer, and 2026's hyperscaler-bundled
prompt-shield products advertise the same goal. The defensive landscape
is converging. The authorisation question is orthogonal and remains
open: even when the model is robust, an agent acting in regulated
environments must produce a portable, auditable record of \emph{what it
was authorised to do, by whom, against what verified policy, and what it
actually did}. VCD's primitives are the first published combination ---
SMT-verified policy completeness, Lean-mechanised attenuation soundness,
content-addressed signed capabilities --- that produces such a record.
The empirical attack-success reduction we report against contemporary
prompt-injection benchmarks is a strong demonstration that the
authorisation layer is also a competitive defence; the durable
contribution is the artefact.

The contribution is neither the SMT verification nor the capability
discipline taken alone --- both have decades of prior art that we
discuss in Section 7 --- but the specific composition, the empirical
demonstration that it holds against state-of-the-art prompt-injection
benchmarks at production-acceptable latency, and the production of a
portable cryptographic audit primitive. We mechanise three soundness
theorems in Lean 4, evaluate against AgentDojo
{[}@debenedetti2024agentdojo{]} and InjecAgent {[}@zhan2024injecagent{]}
alongside contemporary text-side defences (Spotlighting
{[}@hines2024defending{]}, Microsoft Prompt Shields), and release the
reference implementation, Lean development, and benchmark harness as
open source under a permissive licence (Apache-2.0).

Our principal claims are:

\begin{enumerate}
\def\labelenumi{\arabic{enumi}.}
\tightlist
\item
  \textbf{Structural enforcement.} If the gate returns ALLOW for a tool
  call, the call arguments satisfy every constraint in a valid
  descendant of a capability token signed by a pinned issuer; absent
  compromise of the issuer's private key, no agent input ---
  natural-language or otherwise --- can produce an ALLOW for arguments
  that violate the policy.
\item
  \textbf{Provable policy completeness.} Given a tool's JSON Schema in a
  supported fragment, the prover terminates with either a proof that no
  string in the schema's language violates the policy, or a concrete
  counterexample call.
\item
  \textbf{Portable audit primitive.} Every gate decision produces a
  \emph{capability receipt}: a content-addressed, Ed25519-signed record
  citing the issuer's public key, the verified JSON Schema, the cited
  policy proof, the Lean soundness theorem id, the full attenuation
  chain, and the call argument hash. The receipt is verifiable by any
  third party --- including downstream tool implementations, partner
  organisations, and external auditors --- without sharing a secret with
  the issuer. To our knowledge, no prior agent-security work produces an
  artefact with this combination of properties.
\item
  \textbf{Empirical attack-success reduction.} Across three
  frontier-class open-weight models (deepseek-v3.2, deepseek-v4-flash,
  deepseek-v4-pro), three attack families
  (\texttt{important\_instructions}, \texttt{direct},
  \texttt{ignore\_previous}), and four AgentDojo task suites (banking,
  slack, travel, workspace), VCD reduces the tool-call-mediated
  attack-success rate from a no-defence baseline of \textbf{1.4--70.8\%}
  to a benchmark-artefact floor of \textbf{0.0--0.7\%}, while preserving
  \textbf{58.6--91.0\%} benign task completion. The strongest
  contemporary text-side defence we compared (prompt-shields,
  DeBERTa-based PI detector) achieves equivalent ASR but collapses
  benign task completion to \textbf{0.0--36.7\%} on the same cells --- a
  \textbf{+27.9 to +58.6 percentage-point} benign-preservation gap at
  equivalent security. A static upper bound (§6.2.1) establishes that
  the gate's constraint logic denies every catalogued attack across both
  benchmarks' 2,737 user-task × injection-task pairs. Every non-zero
  LLM-driven ASR cell we report (720 banking attempts) maps to the same
  \texttt{user\_task\_15\ ×\ injection\_task\_4} IBAN-coincidence
  benchmark artefact in which the user's authorised tool call
  legitimately matches the attacker's adversarial target; outside that
  single cell, the LLM-driven gate-block rate on attacker-controlled
  tool calls is \textbf{100\%}. Per-call gate latency is well under 100
  microseconds at p50 on commodity hardware, and end-to-end agent wall
  time with VCD enabled is at or below the unprotected baseline on four
  of eight cohorts measured (§6.3.1).
\end{enumerate}

We argue in Section 9 that VCD does not solve prompt injection ---
free-form output attacks and side-channel exfiltration over allowed
parameters are out of scope --- but that it closes the dominant attack
class, and that the trajectory of AI security must be to move
enforcement boundaries out of the model and into structural gates
wherever the grammar permits.

We adopt the convention throughout this paper of being conspicuously
honest about the boundary of our claims. Section 2 enumerates the
out-of-scope attack classes; Section 6.5 reports the configurations in
which VCD did \emph{not} eliminate attack success. The narrowness of the
claim is what makes it credible.

The remainder of this paper is organised as follows. Section 2 fixes the
threat model. Section 3 describes the three primitives and their
composition. Section 4 states and sketches the three soundness theorems.
Section 5 describes the implementation. Section 6 reports the empirical
evaluation. Section 7 places VCD in the related literature. Section 8
discusses limitations.
